\section*{第 \ref{ch:b2:phanerozoic-earth} 章——显生宙的地球}

\begin{solution}{exo:b2:phanerozoic-earth:1}
(a) 巴黎（北纬 \ang{5.1}）与纽约（南纬 \ang{3.6}）；开普敦和悉尼位于 南纬 \ang{60} 附近。
(b) 开普敦：$61.9 - 33.9 = \ang{28.0}$，$28.0 \times 111.2 = \qty{3114}{km}$；纽约：
$40.7 + 3.6 = \ang{44.3}$，\qty{4926}{km}。(c) $3114/300 = \qty{10.4}{km}$ 每百万年，
\qty{1.0}{cm} 每年；$4926/300 = \qty{16.4}{km}$ 每百万年，\qty{1.6}{cm} 每年。
\end{solution}

\begin{solution}{exo:b2:phanerozoic-earth:2}
(a) 72、120.1、50.5 与 \qty{135.4}{Ma}；均值为 $(444 - 66)/4 = \qty{94.5}{Ma}$。(b) 三次事件为：
二叠纪末（西伯利亚，相隔 \qty{1.9}{Ma}）、三叠纪末（CAMP，\qty{0.4}{Ma}）与
白垩纪末（德干，\qty{0.5}{Ma}）。(c) 否：间隔从 50 至
\qty{135}{Ma} 不等；任何一组日期都能算出平均值。
\end{solution}

\begin{solution}{exo:b2:phanerozoic-earth:3}
(a) $\qty{91}{m}/\qty{0.036}{mm.\text{年}^{-1}} = \num{2.5e6}$ 年。(b) $3.7/0.036 = 103$ 倍。(c)
远离任何海岸、水体清澈，来自河流的泥很少能到达海底；
藻类贝壳的沉降是主要沉积物。
\end{solution}

\begin{solution}{exo:b2:phanerozoic-earth:4}
(a) 自河流砂岩至白垩，沉积物向上变深，为\omterm{def:b2:phanerozoic-earth:transgression}{海侵}。
白垩之上的海滩砂记录\omterm{def:b2:phanerozoic-earth:transgression}{海退}。(b) 否：它记录的是相对海平面上升，
而该地沉降或沉积物供应减少也可能造成这一现象。(c)
在其他大陆同一时间找到相同的\omterm{def:b2:phanerozoic-earth:transgression}{海侵}，或
用\omterm{def:b2:basins-compared:backstripping}{回剥}法消除沉降的影响。
\end{solution}

\begin{solution}{exo:b2:phanerozoic-earth:5}
(a) $233/97 = \qty{2.4}{m}$ 每百万年，\qty{2.4}{\micro m} 每年。(b) $3.7/0.0024 =
1540$：今日的上升速度约快 1500 倍。(c) \omterm{def:b2:basins-compared:backstripping}{回剥}曲线：钻探表明
旋回图能准确给出下降的时间，但把幅度夸大了两倍
或更多；新泽西还有自身的偏差（动力地形），因此估计值为 $100 \pm
\qty{50}{m}$。
\end{solution}

\begin{solution}{exo:b2:phanerozoic-earth:6}
(a) $\num{e5} \times \num{e9} \times 2900 = \qty{2.9e17}{kg}$ 的玄武岩，以及 $\num{2.9e17}
\times 0.0075 = \qty{2.2e15}{kg}$ 的二氧化碳，$\num{2.2e15}/0.04401 = \qty{4.9e16}{mol}$。
(b) $\num{278e-6} \times (44.01/28.97) \times \num{5.1e18} = \qty{2.15e15}{kg}$。(c) 一次脉冲
释放的量约相当于 1750 年大气中所含的全部量。(d) $\num{e17}/\num{2.15e15} = 46.5$：接近五十。
\end{solution}

\begin{solution}{exo:b2:phanerozoic-earth:7}
(a) $2630 \times \frac{\pi}{6} \times (\num{e4})^3 = \qty{1.38e15}{kg}$。(b) $\tfrac12 \times
\num{1.38e15} \times (\num{11.2e3})^2 = \qty{8.7e22}{J}$；速度为 \qty{20}{km.s^{-1}} 时，
\qty{2.76e23}{J}。(c) 以 \qty{47}{TW} 的功率计，一年的能量为 $\num{47e12} \times \num{3.156e7} =
\qty{1.48e21}{J}$：相当于 59 年与 186 年。
\end{solution}

\begin{solution}{exo:b2:phanerozoic-earth:8}
(a) 颜色深、纹层细、无潜穴、有机质丰富、含黄铁矿；常伴有
碳同位素的漂移及变暖迹象。(b) 托阿尔期事件对应卡鲁–费勒尔省，OAE~1a 对应
翁通爪哇省。(c) 年代数据的不确定度可达数百万年，且事件与
大火成岩省众多，因此偶然巧合很可能出现；要建立因果联系，需要该省在事件发生前
不久开始活动，并在同一沉积物中要有火山活动的示踪指标，例如汞。
\end{solution}

\begin{solution}{exo:b2:phanerozoic-earth:9}
(a) $60.9 + 5.1 = \ang{66.0}$，$66.0 \times 111.2 = \qty{7339}{km}$，历时 \qty{200}{Ma}：
\qty{36.7}{km} 每百万年，\qty{3.7}{cm} 每年。(b) 超过
自 \qty{300}{Ma} 以来 \qty{1.6}{cm} 每年的两倍。(c) 仅知纬度的变化；
经度变化在古地磁参考系中不受约束，任何绕行也都会增加距离。
\end{solution}

\begin{solution}{exo:b2:phanerozoic-earth:10}
(a) 年龄为 $t$ 的洋底比例为 $p(t) = \frac2T(1 - t/T)$。于是 $\int_0^T tp\,\mathrm{d}t =
\frac2T\bigl(\frac{T^2}2 - \frac{T^2}3\bigr) = \frac T3$，且 $\int_0^T \sqrt t\,p\,\mathrm{d}t =
\frac2T\bigl(\frac23 T^{3/2} - \frac25 T^{3/2}\bigr) = \frac8{15}\sqrt T$。(b) $T = 3 \times 64.2 =
\qty{192.6}{Ma}$；$2500 + \frac8{15} \times 350 \times 13.878 = \qty{5091}{m}$。(c) $T/2 =
\qty{96.3}{Ma}$：\qty{4332}{m}；下降量为 \qty{759}{m}，比
\qty{774}{m} 少 \qty{2}{\percent}。(d) 对于均匀模型：$\frac23\sqrt\tau = \frac23\sqrt{2\bar t} = 0.943\sqrt{\bar t}$；
三角形模型：$\frac8{15}\sqrt{3\bar t} = 0.924\sqrt{\bar t}$。结果主要取决于平均年龄，
远甚于其形状。
\end{solution}

\begin{solution}{exo:b2:phanerozoic-earth:11}
(a) 在 $\Delta s$ 中，只有因子 $1 - f^{-1/2}$ 依赖于 $f$；除以其在 $f = 2$ 时的值，
$1 - 2^{-1/2} = 0.2929$。(b) $f = (1 - 0.2929\,\Delta s/432)^{-2}$：\qty{50}{m} 时为 1.071，
\qty{150}{m} 时为 1.239；比今日快 7 至 \qty{24}{\percent}。(c) $432 \times (1 -
0.8^{-1/2})/0.2929 = \qty{-174}{m}$：海平面将下降 \qty{174}{m}。(d) 洋底年龄分布
仅在老洋底被俯冲并被替换时才改变，此过程所需时间的量级为 $\tau$，
约 \qty{100}{Ma}；海平面则在数千万年的时间尺度上随之变化。
\end{solution}

\begin{solution}{exo:b2:phanerozoic-earth:12}
(a) $2^{-7.4/2} = 0.077$：减少了 \qty{92}{\percent}。(b) $t = 2\log_2(1/0.19) = \qty{4.8}{Ma}$。
(c) 快 $4.8/0.060 = 80$ 倍。(d) 在一个较长的阶内，仅背景灭绝即可
淘汰大部分物种，因此短暂的脉冲叠加在庞大的总数之上，很难将其与之区分；
记录的不完整也会使末现时间向较老一端模糊。
\end{solution}

\begin{solution}{pb:b2:phanerozoic-earth:1}
\textbf{1.} 2500、4250、5304 与 \qty{6000}{m}。
\textbf{2.} $\tau = 2 \times 64.2 = \qty{128.4}{Ma}$；$2500 + 350\sqrt{128.4} = \qty{6466}{m}$。
\textbf{3.} $\num{293e6}/(65000 \times 128.4) = \qty{35.1}{km}$ 每百万年，
\qty{35.1}{mm} 每年：速率很慢。
\textbf{4.} $\frac1\tau\int_0^\tau(2500 + 350\sqrt t)\,\mathrm{d}t = 2500 + \frac23 \times
350\sqrt\tau = \qty{5144}{m}$。
\textbf{5.} 对于超过数千万年的洋底，平方根定律给出的深度偏大，
因此曲线趋于平缓；实际年龄分布并不均匀，年轻洋底最为丰富；并且
沉积物、\omterm{def:b2:oceanic-islands:seamount}{海山}与海台都会抬高洋底。
\textbf{6.} $\num{293e6} \times 5.144 = \qty{1.51e9}{km^3}$，比
整个大洋多 \qty{13}{\percent}：模型偏深。
\textbf{7.} $v = \qty{17.6}{km}$ 每百万年；$17.6 \times 128.4 = \qty{2260}{km}$，
每侧。
\textbf{8.} $t = 1000/17.6 = \qty{56.8}{Ma}$；$2500 + 350\sqrt{56.8} = \qty{5138}{m}$。
\textbf{9.} 在一侧，取 $x = vt$：$\int_0^{v\tau}\bigl(d(\tau) - d(x/v)\bigr)\mathrm{d}x
= vc\int_0^\tau(\sqrt\tau - \sqrt t)\,\mathrm{d}t = \frac13 vc\,\tau^{3/2}$；两侧合计为其两倍。
$\frac23 \times 17.6 \times 0.350 \times 128.4^{3/2} = \frac23 \times 17.6 \times 0.350 \times
1454.9 = \qty{5975}{km^2}$。
\textbf{10.} $65000 \times 5975 = \qty{3.88e8}{km^3}$。
\textbf{11.} $\num{3.88e8}/\num{1.33e9} = 0.29$。
\textbf{12.} 其洋底平均更年轻，因此更热、密度更小、位置更高；沿轴线的
年轻洋底带更宽。
\textbf{13.} $\tau = \qty{64.2}{Ma}$；$2500 + \frac23 \times 350\sqrt{64.2} = \qty{4370}{m}$。
\textbf{14.} $5144 - 4370 = \qty{774}{m}$。
\textbf{15.} $\num{293e6} \times 0.774 = \qty{2.27e8}{km^3}$。
\textbf{16.} $\num{2.27e8}/\num{362e6} = \qty{0.627}{km}$，即 \qty{627}{m}。
\textbf{17.} $\Delta s = 627 \times (3300 - 1025)/3300 = \qty{432}{m}$。
\textbf{18.} $f = \bigl(1 - 0.2929 \times 100/432\bigr)^{-2} = 1.15$：快 \qty{15}{\percent}。
\textbf{19.} $0.149 \times 148 = \qty{22.1}{\mega km^2}$。
\textbf{20.} $48.1 + 0.32 \times 7.7 = \qty{50.6}{\percent}$，\qty{74.9}{\mega km^2}。
\textbf{21.} 水位上升时淹没面积更大，因此同样的体积形成的水层更薄
（见\cref{prop:b2:phanerozoic-earth:flooding}）。
\textbf{22.} 水位在 100、200、300 与 \qty{400}{m} 时的淹没面积为：22.1、41.7、57.3 与
\qty{71.2}{\mega km^2}。各段体积：$(362 + 11.0) \times 0.1 = \qty{37.3e6}{km^3}$，其后依次为 39.4、41.2
与 \qty{42.6e6}{km^3}；累计为 37.3、76.7、117.9 与 \qty{160.5e6}{km^3}。待放置的体积为
\qty{156.4e6}{km^3}，介于 300 与 \qty{400}{m} 之间：$300 + (156.4 - 117.9)/42.6 \times 100
= \qty{390}{m}$。
\textbf{23.} 占今日陆地的 $38.7 + 0.90 \times 9.4 = \qty{47.2}{\percent}$。
\textbf{24.} 将扩张速率加倍将使海平面上升约 \textbf{\qty{390}{m}}，
淹没近半数的今日陆地。
\end{solution}
